\begin{diagram}[
  type=architecture,
  label={fig:sec08-security-layers},
  caption={Security controls sit at five layers along the access path from identity to consumer, not in one undifferentiated "secure the warehouse" bucket; a single vendor tool commonly spans two or three adjacent layers, but each layer's evidence and failure must still be attributable on its own.},
  source={Author's synthesis.},
  description={Five stacked layers, top to bottom, along the access path from identity to consumer: Identity (Authentication, authorization, Secrets); Network (Segmentation, private access); Data (Storage, Classification, Retention, Audit); Application or workload (Processing); Consumers (Access policies). An annotation states that a cloud identity platform commonly spans Identity and Network, and a warehouse policy engine commonly spans Data and Consumers, while audit evidence still has to be attributed to one layer.}
]
\begin{reportarchitecture}[
  annotation={Example span: a cloud identity platform often covers Identity and Network together; a warehouse policy engine often covers Data and Consumers together. Audit evidence still has to be attributed to one layer.},
  annotation-position=top
]
  \layer{Identity}{Authentication, authorization, Secrets}
  \layer{Network}{Segmentation, private access}
  \layer{Data}{Storage, Classification, Retention, Audit}
  \layer{Application or\\ workload}{Processing}
  \layer{Consumers}{Access policies}
\end{reportarchitecture}
\end{diagram}
